% BHU PYQ -- Manually transcribed & error-corrected
% Paper: GLB-601: Palaeontology
% Exam: B.Sc. (Hons.) Semester VI Examination, 2013-14

\documentclass[11pt,a4paper]{article}
\usepackage[a4paper,top=2.5cm,bottom=2.5cm,left=2.8cm,right=2.8cm,headheight=14pt]{geometry}
\usepackage{amsmath,amssymb}
\usepackage[T1]{fontenc}
\usepackage[utf8]{inputenc}
\usepackage{lmodern,microtype,fancyhdr,enumitem,hyperref,lastpage,parskip}

\hypersetup{colorlinks=true,urlcolor=black,linkcolor=black,
  pdftitle={GLB-601 Palaeontology -- BHU B.Sc. Sem VI 2013-14}}
\pagestyle{fancy}\fancyhf{}
\renewcommand{\headrulewidth}{0.4pt}\renewcommand{\footrulewidth}{0.4pt}
\fancyhead[L]{\small   \textit{Banaras Hindu University}}
\fancyhead[R]{\small   \textit{GLB-601: Palaeontology}}
\fancyfoot[C]{\small Page  \thepage\ of \pageref{LastPage}}


\newlist{parts}{enumerate}{2}
\setlist[parts,1]{label=(\alph*),leftmargin=*,itemsep=5pt,topsep=3pt}

\newlist{romanparts}{enumerate}{2}
\setlist[romanparts,1]{label=(\roman*),leftmargin=*,itemsep=5pt,topsep=3pt}

\newcommand{\pts}[1]{\hfill   \textit{[#1]}}


\begin{document}

\begin{center}
  {\large\bfseries BANARAS HINDU UNIVERSITY}\\[2pt]
  {\normalsize B.Sc. (Hons.) --- Semester VI Examination, 2013-14}\\[6pt]
  \rule{0.8\linewidth}{0.5pt}\\[6pt]
  {\large\bfseries Paper: GLB-601 --- Palaeontology}\\[4pt]
  {\normalsize Time: 3 Hours \hfill Maximum Marks: 70}\\[2pt]
  {\small Ref. No.: BS/Sem VI/227}
\end{center}
\rule{\linewidth}{0.4pt}
\smallskip



\noindent   \textit{Instructions: Answer   \textbf{five} questions in all, selecting at least   \textbf{two} each from Sections A and B, and Section-C which is compulsory. All questions carry equal marks.}
\bigskip

\bigskip


\noindent {\large\bfseries SECTION --- A}
\rule{\linewidth}{0.4pt}\smallskip




\noindent   \textbf{Question 1.} Discuss the salient features of Darwin's Theory of Organic Evolution with evidence in support. \pts{14}

\medskip


\noindent   \textbf{Question 2.} Write short notes on   \textbf{any four} of the following: \pts{$4  \times 3\frac{1}{2} = 14$}

\begin{parts}
  \item Index fossils
  \item Development of an individual organism
  \item Stem and cirri of crinoids
  \item Shell coiling in gastropods
  \item Chemical composition of brachiopod shells
\end{parts}

\medskip


\noindent   \textbf{Question 3.} Describe   \textbf{any two} of the following: \pts{$2  \times 7 = 14$}

\begin{parts}
  \item Description and preservation of a fossil
  \item Shell forms in nautiloids
  \item Zooecia of a bryozoan
\end{parts}

\medskip


\noindent   \textbf{Question 4.} Discuss the diagnostic morphological features of bivalves. \pts{14}

\bigskip


\noindent {\large\bfseries SECTION --- B}
\rule{\linewidth}{0.4pt}\smallskip




\noindent   \textbf{Question 5.} Give a brief account of   \textbf{any two} of the following: \pts{$2  \times 7 = 14$}

\begin{parts}
  \item Siwalik vertebrates
  \item Lower Gondwana flora
  \item Derived and living fossils
\end{parts}

\medskip


\noindent   \textbf{Question 6.} Describe the morphological characteristics of ammonoids. \pts{14}

\medskip


\noindent   \textbf{Question 7.} Write short notes on   \textbf{any four} of the following: \pts{$4  \times 3\frac{1}{2} = 14$}

\begin{parts}
  \item Stratigraphic significance of conodonts
  \item Guard in belemnites
  \item Umbilicus and columella in gastropods
  \item Nature of the anterior margin of brachiopods
  \item Types in taxonomy
\end{parts}

\bigskip


\noindent {\large\bfseries SECTION --- C}
\rule{\linewidth}{0.4pt}\smallskip




\noindent   \textbf{Question 8.} Define the following in 2 to 4 sentences: \pts{$7  \times 2 = 14$}

\begin{romanparts}
  \item Protegulum and prodissoconch
  \item Organic-walled microfossils
  \item Operculum
  \item Siliceous microfossils
  \item Ammonitic suture with a neat sketch
  \item Calcareous microfossils
  \item Epirostrum (Epirostum)
\end{romanparts}

\vfill\rule{\linewidth}{0.4pt}

\begin{center}\small
\href{https://bhu-exam.vercel.app/}{Click here to visit our website}\\[4pt]
\href{https://me-aryan.vercel.app/}{Click here to visit our contributor}\\[4pt]
\href{https://quashan-me.vercel.app/}{Click here to visit our contributor}
\end{center}
\end{document}
